% Section 6 of the AAMAS paper, rewritten on 2026-10-06 after Roie's note: no research questions at the
% start; setup, algorithms, results (facts only), then one analysis. Body from the section heading to the
% end of the Analysis paragraph; the appendix zoom captions are changed by build_v6_ors_sec6.py.
% Numbers: the current v6 text and experiments/aamas/section6/out/numbers_kstar.md. Nothing new was run.

\section{Experimental Evaluation}
\label{Sec:experiments}

\begin{figure*}[t]
\centering
\includegraphics[width=.32\textwidth]{s6_sparse_2.pdf}
\includegraphics[width=.32\textwidth]{s6_dense_2.pdf}
\includegraphics[width=.32\textwidth]{s6_scale_free_2.pdf}
\vspace{-8pt}
\caption{Mean solution cost per iteration when solving the random sparse, random dense and scale free problems. MS-SCFG-opt is drawn at iteration 400 and the MGM step of DMS-$k^*$DS-MS-MGM at iterations 1972, 3520 and 1816 (left to right); the time of the search and of MGM is not included.}
\vspace{-8pt}
\label{Fig:random}
\end{figure*}

\paragraph{Setup:}
Following \citet{CohenGZ20}, we evaluate the algorithms on five benchmarks, \emph{Random sparse}, \emph{Random dense}, \emph{Scale-free}, \emph{Graph coloring} and \emph{Meeting scheduling} (their description is in the supplementary material), generating 50 independent instances per benchmark. All algorithms run for 4000 synchronous iterations, and in every iteration, the cost of the assignment selected by the current beliefs is evaluated on the original problem. All damped versions use $\lambda = 0.9$, which was found to be the best in \citet{CohenGZ20}. We report means over the 50 instances; every comparison between two versions is a paired Wilcoxon signed-rank test over the instances, and we call a difference significant when $p < 0.01$.

When we say that an algorithm \emph{converged} we refer to assignment convergence. We declare such convergence if the cost of the algorithm's selected assignment did not change during the last 200 iterations. We report the last iteration at which the cost changed, as the convergence point.

We considered a message from a function-node to a variable-node to be \emph{at a bound} if a single value of the sending variable minimizes the sum of the constraint cost and the incoming belief for every value of the receiving variable. Such a message equals one row of the cost table up to a constant, so its differences are fixed; when the variables are binary this is exactly a message difference that sits at one of the bounds of Section~\ref{sec:split}.

\paragraph{Algorithms:}
The following versions of the algorithm were compared:
\textbf{MS}, standard Min-sum on the original factor graph. \textbf{MS-SCFG}, Min-sum on the symmetric split constraint factor graph, in which every function-node is replaced by two copies holding half of its cost table. \textbf{MS-SCFG-MGM} and \textbf{MS-SCFG-opt} (Section~\ref{Sec:MS-split}): MS-SCFG runs for 400 iterations, the two assignments selected in iterations 398 and 400 define a problem with two values per variable, and this problem is solved by MGM \cite{MaheswaranTBPV04}, or by a centralized branch and bound search over the reduced problem (with a time cap of 300 seconds per instance).
%The search completed on every instance of random sparse, graph coloring and meeting scheduling and on 48 of the 50 scale free instances; on random dense it hit the cap on 48 instances, where we report the best solution found within this time limit.
%On one instance per benchmark the search took 0.4 to 300 seconds, as long as about 250 to 37{,}700 iterations of DMS, and MGM as long as 4 to 9.
\textbf{DMS}, damped Min-sum on the original factor graph. \textbf{DMS-SCFG}, damped Min-sum on the symmetric SCFG. \textbf{DMS-$k$DS} (Section~\ref{Sec:with}): DMS on the original factor graph for $k$ iterations, then the graph is converted to an SCFG as described in Section~\ref{Sec:with}; no $k$ was selected per benchmark, and every $k$ was run on the same 50 instances. \textbf{DMS-$k^*$DS}: DMS-$k$DS in which every instance is split at its own $k^*$, the iteration of the lowest DMS cost within the first 2000 iterations, recorded by the anytime mechanism of \citet{ZivanOP14}. \textbf{DMS-$k^*$DS-MS-MGM}: the same split at $k^*$, after which damping is turned off (MS-split) and MGM selects between the assignments of iterations $t-2$ and $t$, with one $t$ per benchmark: the latest iteration at which an instance starts to repeat its assignments every two iterations (1972, 3520, 1816, 3150 and 3752 on the five benchmarks, in the order above). \textbf{DABP} \cite{DengKL022}, Deep Attentive Belief Propagation, which learns damping factors and message weights on an SCFG with a $0.95/0.05$ split. DABP is centralized, trained, and its iterations require neural nets and a GPU; to plot it on the iteration axis we use the simulated runtime method \cite{SultanikLR08}: with times measured on an NVIDIA RTX 4090, an iteration of DABP takes 6.9, 3.3, 5.6, 6.3 and 8.2 times as long as an iteration of DMS on the five benchmarks, and the DABP curve is stretched by these factors. We include it as a reference for a strong learned centralized version and do not analyze it further. \textbf{DABP-NoSplit} is the same network without the split, stretched by its own factors (3.3, 1.8, 4.3, 4.5 and 6.2).

\begin{figure*}[t]
\centering
\includegraphics[width=.35\textwidth]{s6_graph_coloring_2.pdf}
\includegraphics[width=.35\textwidth]{s6_meeting_scheduling_2.pdf}
\vspace{-8pt}
\caption{Mean solution cost per iteration when solving Graph coloring (left) and Meeting scheduling (right) problems. MS-SCFG-opt is drawn at iteration 400 and the MGM step of DMS-$k^*$DS-MS-MGM at iterations 3150 and 3752 (left and right); the time of the search and of MGM is not included.}
\vspace{-12pt}
\label{Fig:structured}
\end{figure*}

\begin{figure*}[t]
\centering
\includegraphics[width=\textwidth]{s6_bound_fraction.pdf}
\vspace{-8pt}
\caption{Fraction of function-to-variable messages at a bound, per iteration, mean over 50 instances. The dotted vertical line marks the split of DMS-$k$DS ($k=1000$).}
\vspace{-8pt}
\label{Fig:bound}
\end{figure*}

\paragraph{Solution cost:}
Figures~\ref{Fig:random} and~\ref{Fig:structured} present the mean cost per iteration. MS oscillates and produces low quality solutions, and MS-SCFG only slightly improves them. DMS improves the cost slowly. DMS-SCFG ends with a lower mean cost than DMS on all five benchmarks, significantly on random sparse ($-3.6\%$, $p = 6.5\cdot 10^{-5}$), graph coloring ($-75\%$, $p = 2\cdot 10^{-10}$) and meeting scheduling ($-62\%$, $p = 1.2\cdot 10^{-14}$); on random dense its mean is lower by $0.3\%$ while DMS ends lower on 37 of the 50 instances ($p = 0.054$), and on scale free the difference is not significant ($p = 0.25$). The random split of \citet{CohenGZ20}, in which each entry is divided in a ratio drawn uniformly from $[0.4, 0.6)$, is lower than the symmetric split by $0.2$--$0.4\%$ on the three random benchmarks ($p \le 0.006$) and higher by $12\%$ on meeting scheduling ($p = 2.7\cdot 10^{-7}$).

MS-SCFG-MGM and MS-SCFG-opt lower the mean cost of MS-SCFG by $13.3\%$ and $13.5\%$ on random sparse, $6.3\%$ on random dense, $12.2\%$ and $12.5\%$ on scale free, $85\%$ and $88\%$ on graph coloring and $74\%$ and $75\%$ on meeting scheduling, on every instance of the random benchmarks and meeting scheduling and on 49 of the 50 graph coloring instances. Measured against the gap between MS-SCFG and DMS-SCFG, the selection closes $90$--$95\%$ of it on the first four benchmarks and $82$--$84\%$ on meeting scheduling. The search improves on MGM by at most $0.4\%$ on the random benchmarks, $15\%$ on graph coloring and $4.4\%$ on meeting scheduling. DMS-SCFG remains lower than MS-SCFG-opt on every benchmark ($0.7$--$0.9\%$ on the random benchmarks, $p \le 1.6\cdot 10^{-5}$; $59\%$ and $60\%$ on the structured ones).

\begin{table}[t]
\centering
\small
\caption{Mean final cost of DMS-$k$DS relative to DMS-SCFG, in percent (a negative value is a lower cost), 50 instances per benchmark; $^{*}$ marks a significant difference ($p < 0.01$, uncorrected; with Holm's correction over the 26 cells at $0.05$, all marked cells except random sparse at $k = 600$ remain significant); -- marks a value of $k$ that was not run.}
\label{Tab:delay}
\begin{tabular}{lrrrrrr}
$k$ (iterations) & 100 & 200 & 600 & 1000 & 2000 & 3000 \\ \hline
Random sparse & $-0.32^{*}$ & $-0.36^{*}$ & $-0.32^{*}$ & $-0.42^{*}$ & $-0.41^{*}$ & -- \\
Random dense & $-0.22^{*}$ & $-0.31^{*}$ & $-0.32^{*}$ & $-0.33^{*}$ & $-0.36^{*}$ & $-0.40^{*}$ \\
Scale free & $-0.39^{*}$ & $-0.27$ & $-0.59^{*}$ & $-0.65^{*}$ & $-0.75^{*}$ & -- \\
Graph coloring & $+3.5$ & $-10.4$ & $-11.0$ & $-4.1$ & $-8.1$ & -- \\
Meeting sched. & $+0.1$ & $+0.8$ & $-1.4$ & $+1.6$ & $+1.1$ & -- \\
\end{tabular}
\end{table}

Table~\ref{Tab:delay} presents the effect of the moment of the split. On the three random benchmarks every $k$ ends with a lower mean cost than DMS-SCFG, significantly for every $k$ on random sparse and random dense and for every $k$ except 200 on scale free; on random dense the gain grows with $k$, from $0.22\%$ to $0.40\%$, and on scale free from $0.27\%$ at $k = 200$ to $0.75\%$ at $k = 2000$. Within an instance, across the values of $k$ in Table~\ref{Tab:delay}, the Spearman correlation between the cost of DMS at iteration $k$ and the final cost of DMS-$k$DS has a median of $0.58$, $0.66$ and $0.87$ on random sparse, random dense and scale free and is positive on 42, 41 and 40 instances (on 3, 4 and 5 instances every $k$ ends with the same cost). On graph coloring and meeting scheduling no $k$ differs significantly from DMS-SCFG, and the median per-instance correlation is $0.38$ and $0.00$ (on meeting scheduling the mean is within $8.0$--$8.3$ for every $k$).

DMS-$k^*$DS splits every instance at its own $k^*$ (median 613 to 826 iterations on the five benchmarks); up to the split the run equals the DMS run, so, averaged over the instances, the line leaves DMS gradually as one instance after another splits. Since the best iteration of a run is known only in hindsight, we report it as a demonstration of the mechanism at its best starting point rather than as a competing version. On the three random benchmarks the split keeps the best solution of DMS on 30, 35 and 35 of the 50 instances and lowers it on 20, 13 and 15 (on random dense it raises it on 2), and the three lines end at 14{,}507, 99{,}211 and 16{,}591, below every fixed $k$ of Table~\ref{Tab:delay} and below DMS-SCFG ($p \le 9.5\cdot 10^{-13}$). On graph coloring the split changes the best solution on 42 instances (26 lower, 16 higher; mean cost from 36.6 to 28.2, below DMS-SCFG, $p = 4.5\cdot 10^{-4}$) and on meeting scheduling on 41 (37 lower, 4 higher; from 9.25 to 7.97; $p = 0.86$ against DMS-SCFG).

DMS-$k^*$DS-MS-MGM ends at 14{,}511, 99{,}206 and 16{,}590 on the three random benchmarks, not significantly different from DMS-$k^*$DS ($p \ge 0.44$); on 29, 34 and 35 of the 50 instances it ends at exactly the cost of its split point. On graph coloring and meeting scheduling it ends at 60.6 and 13.3, above DMS-$k^*$DS (28.2 and 8.0; $p \le 1.1\cdot 10^{-5}$) and above DMS-SCFG ($p = 0.008$ and $p = 6.4\cdot 10^{-7}$), and below MS-SCFG-MGM on every benchmark ($p \le 4.6\cdot 10^{-7}$).

DABP ends below DMS-SCFG on the three random benchmarks ($0.8\%$, $0.5\%$ and $0.8\%$; $p \le 3.9\cdot 10^{-9}$) and above DMS-$k^*$DS on all three (significantly on scale free, $p = 2.4\cdot 10^{-4}$); on graph coloring and meeting scheduling it ends above DMS-SCFG (39.2 against 34.4, $p = 0.31$, and 28.1 against 8.1, $p = 4.5\cdot 10^{-7}$). Without the split (DABP-NoSplit) its cost is higher on random sparse (15{,}108, $p = 1.4\cdot 10^{-4}$), scale free (17{,}088, $p = 0.004$) and graph coloring (136.4, $p = 3.5\cdot 10^{-10}$); on random dense the difference is not significant ($p = 0.98$), and on meeting scheduling DABP-NoSplit is lower (24.0 against 28.1, $p = 0.36$).

\paragraph{Convergence:}
DMS converged in 17, 33, 22, 8 and 3 of the 50 runs, and among these runs the median convergence iteration was 1088, 966, 717, 758 and 892. DMS-SCFG converged in 49, 50, 49, 40 and 43 runs, with a median convergence iteration of 140, 121, 86, 258 and 208 among them, and DMS-$k$DS ($k = 1000$) converges a median of 46, 32, 50, 118 and 134 iterations after the split. Without damping, only one of the 250 MS-SCFG runs converged: in every benchmark the assignments that MS-SCFG selects in its last 800 iterations switch between two assignments every two iterations (50 of the 50 instances on random sparse, random dense, scale free and meeting scheduling, 48 on graph coloring). The two assignments handed to the selection of Section~\ref{Sec:MS-split}, those of iterations 398 and 400, differ in a median of 36.5, 40, 38 and 41 of the 50 variables on the first four benchmarks (never in fewer than 4, 12 and 7 on the random ones) and in 19 or 20 of the 20 variables on meeting scheduling; they coincide on a single graph coloring instance. After the split of DMS-$k^*$DS-MS-MGM, the undamped run freezes on one assignment on 35, 38 and 41 of the 50 instances of the random benchmarks and alternates between two assignments on 3, 12 and 5 (12, 0 and 4 instances do neither by the end of the run); on graph coloring and meeting scheduling it alternates on 22 and 35 instances and freezes on 11 and 15 (17 graph coloring instances do neither). Among the instances that settle, the median number of iterations after the split until the repetition begins is 4, 4, 4, 62 and 30.

\paragraph{Messages at a bound:}
Figure~\ref{Fig:bound} presents the fraction of messages at a bound. On the original factor graph few messages ever reach a bound: at the end of the run 0.07, 0.47, 0.11, 0.00 and 0.01 of the messages of MS on random sparse, random dense, scale free, graph coloring and meeting scheduling, respectively, and 0.14, 0.81, 0.23, 0.04 and 0.06 of the messages of DMS. On the SCFG of the three random benchmarks most messages reach a bound with or without damping: 0.77, 0.96 and 0.76 for MS-SCFG and 0.77, 0.95 and 0.76 for DMS-SCFG, and DMS-SCFG reaches these levels within the first 200 iterations. On the two structured benchmarks MS-SCFG has only 0.19 and 0.07 of its messages at a bound, DMS-SCFG 0.83 and 0.56. Over the last 800 iterations, the largest change of any variable-to-function message of DMS-SCFG from one step of the implementation (two iterations) to the next is below $10^{-6}$ (maximum norm, each message normalized by subtracting its minimum) in 48, 49, 48, 26 and 37 of the 50 runs, and the selected assignment did not change in 49, 50, 49, 40 and 42 of them. The blue line in Figure~\ref{Fig:bound} is DMS-$k$DS with $k = 1000$. Before the split it follows DMS, with 0.08, 0.61, 0.17, 0.01 and 0.05 of its messages at a bound; 100 iterations after the split 0.75, 0.94, 0.76, 0.66 and 0.56 are at a bound, more than DMS-SCFG has 100 iterations after the start. At every $k$ we tested, 200 iterations after the split the fraction at a bound is within 0.01 of the fraction DMS-SCFG has 200 iterations after the start, or higher. At the split the damping history is cleared, so the first messages after it are not damped; in a control run that clears the history at $k = 1000$ without a split, 100 iterations later 0.08, 0.61, 0.18, 0.01 and 0.05 of the messages are at a bound, as in DMS, and the final cost is not lower than that of DMS.

\paragraph{Analysis:}
\label{Sec:disc}
The measurements tie the two analyses of the paper to the behavior on standard benchmarks. The split brings the messages to the bounds of Section~\ref{sec:split} in large numbers and within the first 200 iterations, on the random benchmarks with or without damping and on the structured ones only with damping; with damping the messages then stop changing. The control run shows that the split, and not the cleared damping history, causes this change. These measurements are consistent with the feedback loop of Section~\ref{sec:split} contributing to the convergence, although they do not establish the hypotheses of its single-cycle theorems in these graphs; the theory covers a single cycle without damping, and the fraction of messages at a bound is our empirical counterpart for graphs with many cycles and larger domains.

Without damping the messages reach the bounds on the random benchmarks, but the assignments alternate between two solutions, as in the lemniscate of Section~\ref{Sec:MS-split}. Selecting between the two candidates of each variable recovers most, but not all, of the gap to the damped version, and the distributed selection keeps almost all of the gain of the search: the two values that each variable holds form a good candidate set, but damping finds a better solution than selecting between them. With damping the alternation disappears and the run converges within a few hundred iterations, where DMS alone needs thousands. This is a statement about speed, not about the cost of the converged solution: on random dense DMS ends lower than DMS-SCFG on 37 of the 50 instances.

The split is not tied to the start of the run: at every iteration we tested it brings the messages to their bounds within the same number of iterations. On the random benchmarks DMS keeps improving its solution long before it converges, so a later split starts from a better solution and ends with a lower cost, and within an instance the cost after the split follows the cost at the split. Where DMS is far from converging, on the structured benchmarks, we did not detect an effect of the moment of the split. Of the Inter-DMS orders of Section~\ref{Sec:with}, the most successful was therefore DMS followed by DMS-SCFG: on the three random benchmarks every fixed $k$ ended below DMS-SCFG, and the split at the best state recorded in the first 2000 iterations, known only in hindsight, ended below every fixed $k$ on all five benchmarks and below DMS-SCFG on graph coloring as well. Turning damping off after that split and selecting with MGM changed nothing on the random benchmarks, where the undamped run freezes within a few iterations of the split, and lost ground on the structured ones, where it moves to worse alternating assignments; damping after the split matters where DMS has not converged. MS-SCFG followed by a search remained behind DMS-SCFG on every benchmark. The advantage DABP draws from its split is not uniform either: it holds on random sparse, scale free and graph coloring, is absent on random dense, and is reversed on meeting scheduling.
